\subsection{CR Alternativi:
CyberRangeCZ}\label{cr-alternativi-cyberrangecz}

\textbf{CyberRangeCZ} e\textquotesingle{} un fork di \textbf{KYPO}.
E\textquotesingle{} una piattaforma basata su OpenStack per fornire
un\textquotesingle ambiente cloud su cui fare il deploy delle macchine.
Utilizza tecnologie moderne come container, IaC (infrastructure as code)
ed i microservizi.

Macchine virtuali, reti, e training sono definiti utilizzando JSON e
YAML. Questi file vengono passati ad Ansible per descrivere il contenuti
della macchina e ad Packer per costruire le immagini delle macchine
virtuali.

Fa uso massivo della virtualizzazione per ricreare fedelmente reti
complesse ed interi sistmi operativi.

La piattaforma si concentra per soddisfare due casi
d\textquotesingle uso principali:

\begin{itemize}
\tightlist
\item
  emulare intere infrastrutture di computer in ambienti controllati
\item
  eseguire più sessioni di training contemporaneamente, gestendo
  centinaia di ambienti virtuali per centinaia o migliaia di
  partecipanti.
\end{itemize}

La piattaforma CyberRangeCZ si occupa di creare e approvvigionare
un\textquotesingle ambiente virtuale emulato costituito da Sandbox e
Pool.

\begin{quote}
\textbf{Sandbox}: \emph{"È un ambiente di test isolato con reti virtuali
che consente agli utenti di connettersi alle macchine virtuali (VM) e
comunicare con il resto della rete, la macchina host e le altre VM.
Sulle VM, l\textquotesingle utente può eseguire programmi, utilizzare
servizi di rete o monitorare il traffico di rete, il tutto senza
influire sull\textquotesingle infrastruttura esterna. La sandbox viene
eseguita all\textquotesingle interno del cloud ed è accessibile da
remoto."} \footnote{\label{fn:ref_13}\url{https://docs.platform.cyberrange.cz/basic-concepts/terminology/\#emulated-virtual-environment}}
\end{quote}

\begin{quote}
\textbf{Pool}: \emph{"Un gruppo di sandbox cloud create a partire dalla
medesima sandbox definition."} \footref{fn:ref_13}
\end{quote}

All\textquotesingle interno di una sandbox e\textquotesingle{} possibile
fornire una \textbf{Topology Definition}, un file che decsrive la
struttura dell\textquotesingle ambiente (host, router, reti, etc.)
mentre il \textbf{Sandbox Provisioning} si occupa di customizzare le
instanze nella topologia.
